\clearpage
\section{Python cho lập trình thi đấu}
\begingroup\small
\lstset{aboveskip=2pt,belowskip=2pt}
\begin{multicols}{2}

\subsection{Nhập và xuất}
\texttt{input()} đọc một dòng và trả \texttt{str}. Khi thay bằng
\texttt{sys.stdin.buffer.readline}, kết quả là \texttt{bytes};
\texttt{int} và \texttt{map(int, ...)} nhận được bytes, còn xử lý
chuỗi Unicode thì phải \texttt{decode()}.
\begin{lstlisting}[language=Python]
import sys
input = sys.stdin.buffer.readline

n = int(input())
a = list(map(int, input().split()))
u, v = map(int, input().split())
s = input().rstrip(b"\r\n").decode()
\end{lstlisting}
\texttt{split()} tách theo mọi khoảng trắng, không phụ thuộc số dấu
cách. Nếu cần giữ khoảng trắng ở hai đầu chuỗi, chỉ bỏ ký tự xuống
dòng như ví dụ; \texttt{strip()} sẽ xóa cả khoảng trắng.
Khi dữ liệu chỉ gồm số và muốn đọc hết một lần:
\begin{lstlisting}[language=Python]
it = iter(map(int, sys.stdin.buffer.read().split()))
n = next(it)
a = [next(it) for _ in range(n)]
\end{lstlisting}
Cách này giữ toàn bộ input trong bộ nhớ. Với dữ liệu lớn, đọc
từng dòng. Nếu số lượng phần tử không đủ, \texttt{next(it)}
phát sinh \texttt{StopIteration}.

\begin{lstlisting}[language=Python]
print(x, y)                 # cach nhau mot dau cach
print(*a)                   # in cac phan tu tren mot dong
out = [str(x) for x in answers]
sys.stdout.write("\n".join(out))
\end{lstlisting}
Gom kết quả rồi \texttt{join} hữu ích khi có rất nhiều dòng;
\texttt{join} cần các phần tử là \texttt{str}.

\subsection{Vòng lặp, dãy và chuỗi}
\texttt{range(r)} chạy từ $0$ đến $r-1$;
\texttt{range(l,r)} không chứa $r$. \texttt{enumerate} giúp lấy cả
chỉ số và giá trị; đặt \texttt{start=1} khi đề đánh số từ $1$.
\begin{lstlisting}[language=Python]
for i in range(n):
    pass
for i in range(n - 1, -1, -1):
    pass
for i, x in enumerate(a, start=1):
    pass

a = [0] * n
grid = [[0] * m for _ in range(n)]
part = a[l:r]               # chi so l,...,r-1
rev = a[::-1]
\end{lstlisting}
Không tạo ma trận bằng \texttt{[[0] * m] * n}: các hàng khi đó
trỏ tới cùng một list. \texttt{a[l:r]} tạo list mới.
\begin{lstlisting}[language=Python]
words = "a b c".split()
s = " ".join(words)
letters = list(s)
code = ord('A')             # 65
ch = chr(65)                # 'A'
\end{lstlisting}

\subsection{Sắp xếp và tìm vị trí}
\texttt{sort()} sửa list tại chỗ và trả \texttt{None};
\texttt{sorted()} trả list mới. Python sắp xếp ổn định:
những phần tử có khóa bằng nhau giữ nguyên thứ tự cũ.
\begin{lstlisting}[language=Python]
a.sort()
a.sort(reverse=True)
pairs.sort(key=lambda p: (p[1], -p[0]))
b = sorted(a)

from bisect import bisect_left, bisect_right
i = bisect_left(a, x)       # dau tien >= x
j = bisect_right(a, x)      # dau tien > x
count_x = j - i
\end{lstlisting}
Hai hàm \texttt{bisect} yêu cầu đoạn đang xét đã sắp xếp theo
cùng quy tắc; trả vị trí chèn trong $[0,n]$.

\subsection{Từ điển, tập hợp, hàng đợi và heap}
\begin{lstlisting}[language=Python]
from collections import Counter, defaultdict, deque
from heapq import heappush, heappop

freq = Counter(a)
cnt = freq[x]                # 0 neu x chua xuat hien
g = defaultdict(list)
g[u].append(v)
seen = set()
seen.add(u)
q = deque([u])
q.append(v)
first = q.popleft()
heap = []                    # min-heap
heappush(heap, (dist, v))
dist, v = heappop(heap)
\end{lstlisting}
Với \texttt{dict} thường, dùng \texttt{d.get(key, 0)} để nhận
giá trị mặc định mà không tạo khóa mới. \texttt{defaultdict(list)}
tạo list rỗng khi truy cập khóa chưa có. Heap Python là min-heap;
muốn lấy số lớn nhất có thể đẩy \texttt{-x} rồi đổi dấu khi lấy ra.

\subsection{Các lỗi cú pháp dễ gặp}
\begin{itemize}
\setlength{\itemsep}{2pt}
\item Python dùng thụt lề để xác định khối lệnh; điều kiện và
      vòng lặp kết thúc phần đầu bằng dấu \texttt{:}.
\item \texttt{/} là chia số thực; \texttt{//} là chia lấy sàn,
      khác phép cắt về $0$ khi số âm. \texttt{**} là lũy thừa,
      \verb|^| là XOR, không phải lũy thừa.
\item \texttt{map(...)} và \texttt{zip(...)} trả iterator;
      chuyển thành \texttt{list(...)} nếu cần truy cập nhiều lần.
\item Kiểm tra giới hạn thời gian trước khi dùng Python cho vòng lặp
      rất lớn hoặc thuật toán đệ quy sâu. Số nguyên lớn, \texttt{pow}
      modulo và các hàm \texttt{math} xem mục 12.5.
\end{itemize}

\end{multicols}
\endgroup
